\documentclass[11pt, a4paper]{article}

% ── Packages ──────────────────────────────────────────────────────────────────
\usepackage[margin=1in]{geometry}
\usepackage[T1]{fontenc}
\usepackage[utf8]{inputenc}
\usepackage{lmodern}
\usepackage{microtype}
\usepackage{graphicx}
\usepackage{hyperref}
\usepackage[dvipsnames]{xcolor}
\usepackage{enumitem}
\usepackage{titlesec}
\usepackage{longtable}
\usepackage{tabularx}
\usepackage{booktabs}
\usepackage{listings}
\usepackage{fancyhdr}
\usepackage{array}
\usepackage{float}
\usepackage{pgfgantt}
\usepackage{caption}
\usepackage{parskip}
\usepackage{wrapfig}

% ── Text overflow fix ─────────────────────────────────────────────────────────
\emergencystretch=5em

% ── HPX / STE||AR brand palette (extracted from HPX_STELLAR_blue.png) ─────────
% Dominant logo colour: deep navy #0A145A
\definecolor{hpxnavy}{RGB}{10,20,90}       % logo primary — section rules, headers
\definecolor{hpxmid}{RGB}{0,70,180}        % mid-blue — hyperlinks
\definecolor{hpxfaint}{RGB}{232,237,248}   % pale brand blue — title background
% Alias so existing \color{linkblue} usage stays consistent
\colorlet{linkblue}{hpxmid}
\definecolor{codebg}{RGB}{245,245,245}
\definecolor{codecomment}{RGB}{100,100,100}
\definecolor{merged}{RGB}{30,130,60}
\definecolor{openpr}{RGB}{200,90,0}
\definecolor{rowgray}{RGB}{240,240,240}
\definecolor{bugred}{RGB}{180,30,30}
\definecolor{diffrem}{RGB}{180,30,30}
\definecolor{diffadd}{RGB}{30,130,60}

% ── Hyperref ──────────────────────────────────────────────────────────────────
\hypersetup{
  colorlinks=true, linkcolor=linkblue, urlcolor=linkblue, citecolor=linkblue,
  pdftitle={GSoC 2026 -- Add HPX to Compiler Explorer},
  pdfauthor={Jatin Sharma},
}

% ── Section formatting ────────────────────────────────────────────────────────
\titleformat{\section}{\Large\bfseries\color{hpxnavy}}{}{0em}{}{\color{hpxnavy}\hrule height 1.2pt}
\titleformat{\subsection}{\large\bfseries\color{hpxnavy}}{}{0em}{}
\titleformat{\subsubsection}{\normalsize\bfseries\color{hpxnavy}}{}{0em}{}
\titlespacing*{\section}{0pt}{14pt}{6pt}
\titlespacing*{\subsection}{0pt}{10pt}{4pt}

% ── Header / Footer ──────────────────────────────────────────────────────────
\pagestyle{fancy}
\fancyhf{}
\lhead{\small\color{hpxnavy}\bfseries GSoC 2026 | Add HPX to Compiler Explorer}
\rhead{}
\cfoot{\small\color{hpxnavy}\thepage}
\renewcommand{\headrulewidth}{1pt}
\renewcommand{\headrule}{\hbox to\headwidth{\color{hpxnavy}\leaders\hrule height \headrulewidth\hfill}}

% ── Code listings ─────────────────────────────────────────────────────────────
\lstdefinestyle{code}{
  backgroundcolor=\color{codebg}, basicstyle=\ttfamily\small,
  keywordstyle=\bfseries, commentstyle=\color{codecomment}\itshape,
  breaklines=true, frame=single, framerule=0.4pt,
  rulecolor=\color{black!30}, numbers=left,
  numberstyle=\tiny\color{gray}, numbersep=5pt,
  xleftmargin=14pt, aboveskip=8pt, belowskip=8pt,
  showstringspaces=false,
}
\lstdefinestyle{diff}{
  backgroundcolor=\color{codebg}, basicstyle=\ttfamily\small,
  breaklines=true, frame=single, framerule=0.4pt,
  rulecolor=\color{black!30},
  xleftmargin=6pt, aboveskip=8pt, belowskip=8pt,
  showstringspaces=false,
}
\lstset{style=code}

% ── Helpers ───────────────────────────────────────────────────────────────────
\newcommand{\code}[1]{\texttt{#1}}
\newcommand{\pr}[1]{\href{https://github.com/STEllAR-GROUP/hpx/pull/#1}{\texttt{\##1}}}
\newcommand{\cefile}[1]{\texttt{#1}}

% ══════════════════════════════════════════════════════════════════════════════
\begin{document}

% ── Title page ───────────────────────────────────────────────────────────────
{\color{hpxnavy}\rule{\linewidth}{2.5pt}}
\vspace{2pt}
\begin{center}
  % HPX + STE||AR Group combined logo
  
\includegraphics[height=2.0cm,keepaspectratio]{hpx_stellar_logo}\\[10pt]
  {\color{hpxnavy}\rule{0.55\linewidth}{0.6pt}}\\[8pt]
  {\LARGE\bfseries\color{hpxnavy} Google Summer of Code 2026}\\[4pt]
  {\Large\bfseries\color{black} Add HPX to Compiler Explorer (godbolt.org)}\\[6pt]
  {\normalsize\color{hpxnavy}\textsc{STEllAR-GROUP / HPX}}\\[4pt]
  {\normalsize Jatin Sharma \quad{\color{hpxnavy}$\cdot$}\quad Medium project (175~hours)}
\end{center}
\vspace{2pt}
{\color{hpxnavy}\rule{\linewidth}{2.5pt}}
\vspace{10pt}

% ══════════════════════════════════════════════════════════════════════════════
\section*{Abstract}

Compiler Explorer (\url{https://godbolt.org/}) is a widely popular web-based
application which provides easy access to multiple C++ compilers and
environments, allowing users to write, test, and share anything from small
C++ scripts to large CMake-based projects quickly.  Given its versatility, we
believe it would be convenient for HPX to have an integration with Compiler
Explorer.  A preliminary idea is that we will maintain our own fork of
Compiler Explorer (which is open source) and experiment with the HPX
integration locally before making the integration public through a Pull
Request to Compiler Explorer.  The result could even be constrained to an
environment somewhat similar to Compiler Explorer just for HPX, where
prospective users would experiment with quick HPX scripts before downloading,
building, and running the whole library.

I have already \textbf{built a working proof of concept}
(HPX~2.0.0 compiling and executing inside a local CE instance) and
\textbf{published the fix} at
\href{https://github.com/GitMasterJatin/Integrating-HPX-with-Compiler-Explorer-}%
{github.com/GitMasterJatin/Integrating-HPX-with-Compiler-Explorer-}.
This proposal converts that validated prototype into three
upstream Pull Requests across the HPX and CE repositories, giving 3.5~million
users the ability to try HPX parallelism in seconds, replacing a 30-minute
local build with a single checkbox.

% ══════════════════════════════════════════════════════════════════════════════
\section*{Personal Details}

\begin{wrapfigure}{r}{3.5cm}
  \vspace{-20pt}
  
\includegraphics[width=3.3cm]{jatin_photo}
  \vspace{-15pt}
\end{wrapfigure}

\begin{tabular}{@{}ll}
  \textbf{Name}        & Jatin Sharma \\
  \textbf{University}  & Indian Institute of Technology (BHU), Varanasi \\
  \textbf{Programme}   & Integrated Dual Degree (Graduating 2027) \\
  \textbf{Email}       & \href{mailto:jatin.sharma.cd.civ22@itbhu.ac.in}%
                               {jatin.sharma.cd.civ22@itbhu.ac.in} \\
  \textbf{GitHub}      & \href{https://github.com/GitMasterJatin}%
                               {github.com/GitMasterJatin} \\
  \textbf{LinkedIn}    & \href{https://www.linkedin.com/in/jatin-sharma-70a26825b/}%
                               {linkedin.com/in/jatin-sharma-70a26825b} \\
  \textbf{Website}     & \href{https://www.sharmajatin.com}{www.sharmajatin.com} \\
  \textbf{Phone}       & (+91)\,9996499472 \\
  \textbf{Project}     & Medium, 175 hours \\
  \textbf{Timezone}    & IST (UTC+5:30) \\
\end{tabular}

% ══════════════════════════════════════════════════════════════════════════════
\section*{Background Information}

\subsection*{Education}

Integrated Dual Degree at \textbf{IIT~(BHU), Varanasi} (graduating 2027).
Core coursework directly relevant to this project:
\begin{itemize}[leftmargin=2em, itemsep=1pt]
  \item \textbf{Data Structures \& Algorithms:} complexity analysis,
        graph algorithms, dynamic programming.
  \item \textbf{Operating Systems:} process scheduling, memory management,
        IPC, cgroups, namespaces (relevant to \code{nsjail}).
  \item \textbf{Computer Architecture:} cache hierarchies, NUMA topology,
        instruction-level parallelism (relevant to HPX's hwloc integration).
  \item \textbf{Parallel Computing:} shared-memory models, task
        decomposition, synchronisation primitives (relevant to HPX's
        work-stealing scheduler and thread pools).
\end{itemize}

\subsection*{Programming Background}

\textbf{Software Development Intern, Fiverse}
\begin{itemize}[leftmargin=2em, itemsep=2pt]
  \item Built a full-stack investor--company matching platform with Nginx
        reverse proxy and 15+ REST routes in Node.js/Express, the same
        stack powering CE's \cefile{base-compiler.ts} and
        \cefile{options-handler.ts}.
\end{itemize}

\textbf{Order Execution \& Management System (C++)}
\begin{itemize}[leftmargin=2em, itemsep=2pt]
  \item Low-latency order placement with multithreading, async HTTP,
        and thread-safe data structures. Complex C++ library integration
        under tight runtime constraints that mirrors HPX's behaviour inside
        CE's 10-second \code{nsjail} sandbox.
\end{itemize}

\textbf{SyncCanvas: Real-Time Collaborative Drawing}
\begin{itemize}[leftmargin=2em, itemsep=2pt]
  \item WebSocket-based multi-user sessions in a TypeScript monorepo
        (Turborepo/pnpm), a structure that mirrors CE's own monorepo.
        CI/CD pipeline experience maps directly to CE's nightly library
        builds.
\end{itemize}

\subsection*{Interest in HPX}

HPX implements the complete C++ standard parallel algorithm suite
(\code{hpx::for\_each}, \code{hpx::transform\_reduce}, \code{hpx::sort})
and extends it with futures, continuations, and a work-stealing runtime.
Code written with \code{hpx::execution::par} is structurally identical to
\code{std::execution::par}, making HPX the ideal teaching and prototyping tool
for modern C++ parallelism.

My motivation is concrete: every time I wanted to show a classmate that
\code{hpx::transform\_reduce(par, ...)} outperforms the sequential version,
I had to walk them through a 30-minute build instead of sharing a godbolt link.
To prove the fix is feasible, I built HPX~2.0.0 inside a local CE instance
and validated all three workflows: assembly inspection, compilation, and
execution (Figures~\ref{fig:ce-lib}--\ref{fig:ce-asm}).
Under the mentorship of \textbf{Hartmut Kaiser} and \textbf{Giannis
Gonidelis}, this integration will follow STE\texttt{||}AR Group's
contribution standards and CE's library addition guidelines.

\subsection*{Contributions to HPX}

Ten merged PRs and one open PR across bug fixes, CMake modernisation,
scheduler performance, and API cleanup:

\begin{table}[!htbp]
\centering\small
\begin{tabular}{@{}rp{6cm}cp{5cm}@{}}
\toprule
\textbf{\#} & \textbf{Pull Request} & \textbf{Status}
  & \textbf{Skill $\to$ This Project} \\
\midrule
\rowcolor{rowgray}
1  & Fix heap corruption in \code{partial\_sort} filter \pr{7051}
   & \textcolor{merged}{\bfseries M}
   & Valgrind/ASan diagnostics $\to$ debugging compiler-specific build failures \\
2  & cmake: remove \code{CXX20\_STD\_RANGES\_ITER\_SWAP} test \pr{7032}
   & \textcolor{merged}{\bfseries M}
   & Obsolete feature-test removal $\to$ writing \code{skip\_compilers} regex \\
\rowcolor{rowgray}
3  & Fix numa allocator test \pr{7027}
   & \textcolor{merged}{\bfseries M}
   & NUMA/hwloc understanding $\to$ hwloc fallback strategy on CE nodes \\
4  & Fix heap corruption in \code{sort\_thread} partitioning \pr{7018}
   & \textcolor{merged}{\bfseries M}
   & Multi-threaded correctness $\to$ validating HPX thread pool in sandbox \\
\rowcolor{rowgray}
5  & cmake: remove \code{CXX20\_SOURCE\_LOCATION} test \pr{6996}
   & \textcolor{merged}{\bfseries M}
   & CMake feature detection $\to$ preset design for \code{godbolt-minimal} \\
6  & cmake: remove \code{CXX20\_STD\_DEFAULT\_SENTINEL} test \pr{6977}
   & \textcolor{merged}{\bfseries M}
   & HPX CI matrix knowledge $\to$ compiler compatibility matrix \\
\rowcolor{rowgray}
7  & cmake: remove \code{CXX20\_STD\_BIT\_CAST} test \pr{6976}
   & \textcolor{merged}{\bfseries M}
   & CMake cache variables $\to$ \code{libraries.yaml} flag design \\
8  & Convert \code{message\_buffer\_append\_state} to enum class \pr{6923}
   & \textcolor{merged}{\bfseries M}
   & Code modernisation $\to$ C++20-clean example code for CE \\
\rowcolor{rowgray}
9  & Add subtraction tests for \code{gid\_type} \pr{6913}
   & \textcolor{merged}{\bfseries M}
   & Unit testing distributed primitives $\to$ E2E pipeline testing \\
10 & Removed incomplete line in README \pr{6898}
   & \textcolor{merged}{\bfseries M}
   & Documentation quality $\to$ maintenance guide + blog post \\
\rowcolor{rowgray}
11 & Improve SPSC queue in work-requesting scheduler \pr{7114}
   & \textcolor{openpr}{\bfseries O}
   & Scheduler internals $\to$ understanding runtime init overhead \\
\bottomrule
\end{tabular}
\caption{HPX contributions. M = Merged, O = Open. Right column maps each
  PR's demonstrated skill to a specific task in this project.}
\end{table}

\subsection*{Open-Source Projects}

\href{https://github.com/GitMasterJatin/Integrating-HPX-with-Compiler-Explorer-}%
{\textbf{HPX $\times$ Compiler Explorer: Proof of Concept}} | local CE
instance with HPX~2.0.0 integration, build scripts, example programs, and
API smoke tests. \textbf{Directly demonstrates this proposal's feasibility.}

\href{https://github.com/GitMasterJatin/Parallel-matrix-multiplication-implemented-in-HPX}%
{Parallel matrix multiplication in HPX}, includes CMake build file as
required by the STE\texttt{||}AR submission template.

\subsection*{Technical Skills and Development Environment}

I possess advanced skills in:
\begin{itemize}[leftmargin=2em, itemsep=4pt]
  \item \textbf{C++:} 3+ years with Modern C++ (C++17/20), including lambdas,
        smart pointers, templates, concurrency, and template metaprogramming.

  \item \textbf{Python:} 3+ years of experience in scientific computing,
        scripting, CPython internals, and async/await.

  \item \textbf{Parallel Computing:} Experience with multithreading
        (\code{pthreads}, \code{std::thread}) and MPI, which complements
        HPX's task-based asynchronous model.

  \item \textbf{Boost Libraries:} Proficient with Boost.Asio, Boost.Thread,
        and related components.

  \item \textbf{Tools:} Skilled in Git/GitHub workflows, CMake (with
        FetchContent/ExternalProject), Linux, and debugging (gdb, Valgrind).
\end{itemize}

My environment features VS~Code with clangd and Cursor for AI-assisted
coding.  I maintain a reproducible setup with HPX built from source.

% ══════════════════════════════════════════════════════════════════════════════
\section*{Project Proposal}

% ── Problem ───────────────────────────────────────────────────────────────────
\subsection*{Problem Description}

Compiler Explorer (\url{https://godbolt.org/}) serves 3.5~million monthly
users with access to 40+ C++ compilers and 130+ libraries.  It is the de
facto tool for sharing C++ snippets, comparing compiler output, and
benchmarking library performance.  Comparable parallelism libraries enjoy
first-class CE support:

\begin{table}[H]
\centering\small
\begin{tabularx}{\linewidth}{@{}lccX@{}}
\toprule
\textbf{Library} & \textbf{CE Status} & \textbf{Versions}
  & \textbf{Build Complexity} \\
\midrule
\rowcolor{rowgray}
Kokkos   & \textcolor{merged}{Working} & 4.0--5.0 (multiple versions)
  & Static, CMake, FetchContent deps \\
Catch2   & \textcolor{merged}{Working} & 3.x, 2.x & Header-only \\
\rowcolor{rowgray}
fmt      & \textcolor{merged}{Working} & 10.x--8.x & Light CMake \\
TBB      & \textcolor{merged}{Working} & 2021.x  & System install \\
\rowcolor{rowgray}
\textbf{HPX} & \textcolor{red}{\bfseries Broken} & 1.11.0 only
  & \textbf{3 deps, custom CMake, dynamic linking} \\
\bottomrule
\end{tabularx}
% Kokkos is the closest analogue: multiple working versions with static linking.
  % [VERIFY: check current Kokkos entry in CE before submission]
\end{table}

Today, trying HPX requires cloning the repository, installing Boost ($\geq$1.71),
hwloc, and Asio, configuring CMake, and building for 10--30 minutes.  The
barrier is enormous for a newcomer who wants to see HPX parallelism in action.

% ── Root Cause Analysis ──────────────────────────────────────────────────────
\subsection*{Root Cause Analysis: Why HPX Is Broken on CE}

An HPX entry already exists in CE's production configuration.  After auditing
\cefile{c++.amazon.properties} (lines~5759--5767) and tracing through CE's
build pipeline (\cefile{library\_builder.py}), I identified \textbf{six root
causes}.  The current broken entry:

% [GROUNDED: compiler-explorer/etc/config/c++.amazon.properties:5759-5767]
\begin{lstlisting}[style=diff, caption={Current broken HPX entry (c++.amazon.properties:5759--5767)}, label=lst:broken]
libs.hpx.name=HPX
libs.hpx.versions=v1110
libs.hpx.url=https://github.com/STEllAR-GROUP/hpx
libs.hpx.description=The C++ Standard Library for Concurrency and Parallelism
libs.hpx.packagedheaders=true
libs.hpx.versions.v1110.version=1.11.0
libs.hpx.versions.v1110.path=/app/hpx/hwloc_installed/include:/opt/compiler-explorer/libs/boost_1_84_0
libs.hpx.versions.v1110.libpath=/app/hpx/hwloc_installed/lib
libs.hpx.versions.v1110.liblink=hpx_wrap:hpx_init:hpx:hpx_core:hwloc:dl:rt
\end{lstlisting}

\subsubsection*{Bug \#1: Hardcoded Boost Path}

\textbf{What users see:} \code{fatal error: boost/version.hpp: No such file or directory}.

\textbf{Root cause:} Line~5765 references
\code{/opt/compiler-explorer/libs/boost\_1\_84\_0}.  When CE updated to
Boost~1.86.0, the path became stale.  The entry should use \code{\%DEP0\%}
variable interpolation. CE's \cefile{library\_builder.py} calls
\code{expand\_make\_arg()} to replace \code{\%DEP0\%} with the resolved
Boost path at build time.

% [GROUNDED: cmake/HPX_SetupBoost.cmake:39 — Boost_MINIMUM_VERSION = 1.71]
\textbf{Fix:} In \code{libraries.yaml}: add
\code{depends: [libraries/c++/boost\_bin 1.86.0]} and use
\code{-DBOOST\_ROOT=\%DEP0\%}.  HPX requires Boost~$\geq$1.71
(\cefile{cmake/HPX\_SetupBoost.cmake:39}); CE ships 1.86.0.

\subsubsection*{Bug \#2: Dynamic Linking in a Static-Only Sandbox}

\textbf{What users see:}
\code{error while loading shared libraries: libhpx\_core.so: cannot open}.

\textbf{Root cause:} The entry uses \code{liblink} (shared linking).
CE's \code{nsjail} sandbox (\cefile{etc/nsjail/user-execution.cfg}) mounts
\code{/opt/compiler-explorer} read-only and does not augment
\code{LD\_LIBRARY\_PATH} for HPX's \code{.so} files.  All successfully
integrated compiled libraries (Kokkos, Abseil, re2) use \code{staticliblink}.

\textbf{Fix:} \code{lib\_type: static} in YAML; \code{staticliblink} in
properties; \code{-DHPX\_WITH\_STATIC\_LINKING=ON} in CMake args.

\subsubsection*{Bug \#3: Only v1.11.0, Outdated Target}

\textbf{What users see:} Even if linking worked, v1.11.0 lacks the
\code{godbolt-minimal} CMake preset and \cefile{sandbox.hpp}.

% [GROUNDED: CMakeLists.txt:69-71 — HPX_VERSION = 2.0.0]
\textbf{Fix:} Target v2.0.0 (current release, per \cefile{CMakeLists.txt:69--71})
as primary; retain v1.11.0 as fallback.

\subsubsection*{Bug \#4: No \code{skip\_compilers}}

\textbf{What users see:} Nightly builds attempted on GCC~4.x / Clang~3.x
(no C++17, let alone C++20), wasting build capacity.

% [GROUNDED: .github/workflows/ — CI uses Clang 20, GCC 11+]
\textbf{Fix:} \code{skip\_compilers} listing all CE compiler IDs for
GCC~$\leq$11 and Clang~$\leq$14, e.g.\ \code{g101}, \code{g111},
\code{clang1400}.  CE's build system matches IDs exactly (not regex).
Derived from HPX's CI matrix (Clang~20, GCC~12+ for reliable C++20).
The full list will be refined in W2 against CE's live compiler inventory.

\subsubsection*{Bug \#5: No Default Example}

\textbf{What users see:} Empty editor.  No indication that
\code{hpx\_main.hpp} is required for runtime initialisation.

\textbf{Fix:} \code{example=hpx-sandbox} in properties, pointing to a
\code{transform\_reduce} example with \code{sandbox.hpp} benchmarking.

\subsubsection*{Bug \#6: \code{sandbox.hpp} Undiscoverable}

% [GROUNDED: libs/full/include/include/hpx/experimental/sandbox.hpp]
\textbf{What users see:} HPX ships a CE-specific header
(\cefile{libs/full/include/include/hpx/experimental/sandbox.hpp}%
\footnote{The doubled \code{include/include} is HPX's actual directory
structure: the outer \code{include/} is the CMake install-tree source
and the inner \code{include/} mirrors the public header prefix.}) with
environment introspection, benchmarking, and scaling telemetry, but
no one knows it exists.

\textbf{Fix:} The default example (Bug~\#5) will \code{\#include
<hpx/experimental/sandbox.hpp>}.

% ── Proposed Solution ─────────────────────────────────────────────────────────
\subsection*{Proposed Solution}

Following the project description, \textit{``we will maintain our own fork of
Compiler Explorer and experiment with the HPX integration locally before making
the integration public through a Pull Request''}, and under the guidance of
mentors \textbf{Hartmut Kaiser} and \textbf{Giannis Gonidelis}, I propose a
\textbf{three-repository, phased integration}:

\begin{description}[leftmargin=1.5em, itemsep=3pt]
  \item[PR\,\#1: \code{STEllAR-GROUP/hpx}]
    Refine the \code{godbolt-minimal} CMake preset (already in
    \cefile{CMakePresets.json}), improve \code{sandbox.hpp}, add integration
    documentation (\cefile{CONTRIBUTING\_CE.md}).

  \item[PR\,\#2: \code{compiler-explorer/infra}]
    Fix \code{libraries.yaml}: static linking, \code{\%DEP0\%} for Boost,
    v2.0.0 target, \code{skip\_compilers}, Ninja, \code{FETCH\_ASIO}.

  \item[PR\,\#3: \code{compiler-explorer/compiler-explorer}]
    Add \code{c++.amazon.properties} entries with correct metadata, link
    flags, and default example.
\end{description}

\textbf{Fork strategy.}
All development happens on a personal fork of \code{compiler-explorer/infra}
and \code{compiler-explorer/compiler-explorer}; the fork carries only the
HPX-specific YAML entry, properties additions, and default example; no
modifications to CE core code.
The fork is rebased onto CE \code{main} weekly (or before each draft PR
update) to prevent drift.
If upstream review extends beyond GSoC, the fork remains a
self-contained, deployable fallback that STE\texttt{||}AR Group can host
as an HPX-branded instance with a single \code{git clone}.

Each PR is self-contained.  Draft PRs open in Week~5 for early CE maintainer
feedback, providing 5+ weeks of review buffer before final evaluation.

\textbf{Upstream PR vs.\ HPX-branded instance.}
The project description mentions two approaches: integrating HPX into
the public CE instance and maintaining an HPX-branded fork.  The primary
deliverable is the upstream PR path (PR\,\#1--\#3 above), because it
reaches CE's full 3.5~million user base with zero maintenance overhead
for STE\texttt{||}AR Group.  The HPX-branded instance, with extended
timeouts, curated examples, and potential task-graph visualisation, is
a stretch goal (Section~``Stretch goals'').  If upstream PRs stall in
review, the fork doubles as a working fallback that STE\texttt{||}AR
Group can host independently.

\subsection*{Deliverables with Acceptance Criteria}

\begin{table}[H]
\centering\small
\begin{tabularx}{\linewidth}{@{}clXl@{}}
\toprule
\textbf{ID} & \textbf{Deliverable} & \textbf{Acceptance Criterion}
  & \textbf{Repo} \\
\midrule
\rowcolor{rowgray}
D1 & Fixed \code{libraries.yaml} &
  HPX v2.0.0 builds on GCC~\{12,14\} and Clang~\{17,20\} in $<$600\,s;
  \code{conan export-pkg} produces valid static \code{.a} files.
  & CE infra \\
D2 & Properties file &
  \code{c++.amazon.properties} entries pass CE's property linter;
  HPX appears in library dropdown with correct metadata.
  & CE main \\
\rowcolor{rowgray}
D3 & Default example &
  Example auto-loads when HPX is selected; compiles on all
  supported compilers; execution output includes
  \code{sandbox::benchmark} results.
  & CE main \\
D4 & Updated CMake preset &
  \code{cmake --preset godbolt-minimal} succeeds with GCC~14 and
  Clang~20; produces \code{libhpx.a}, \code{libhpx\_core.a},
  \code{libhpx\_init.a}, \code{libhpx\_wrap.a}.
  File: \cefile{CMakePresets.json}.
  & HPX \\
\rowcolor{rowgray}
D5 & Compiler matrix &
  Spreadsheet covering GCC~10--15 and Clang~15--20 with build
  status, binary size, and build time for each.
  & (PR attach) \\
D6 & Three upstream PRs &
  PR\,\#1 merged into HPX; PR\,\#2 and PR\,\#3 submitted to CE
  with build logs for $\geq$8 compilers.
  & all three \\
\rowcolor{rowgray}
D7 & Documentation &
  Maintenance guide (\cefile{CONTRIBUTING\_CE.md} in HPX),
  blog post published, CE integration walkthrough.
  & HPX \\
\bottomrule
\end{tabularx}
\caption{Deliverables with falsifiable acceptance criteria and named
  upstream repositories.}
\end{table}

% ── Proof of Concept ──────────────────────────────────────────────────────────
\subsection*{Proof of Concept: Local Validation}

To demonstrate feasibility \textbf{before writing this proposal}, I built and
validated a complete local integration.  The PoC is published at:
\begin{center}
\url{https://github.com/GitMasterJatin/Integrating-HPX-with-Compiler-Explorer-}
\end{center}

\subsubsection*{Setup}

HPX~2.0.0 was built from source using the \code{godbolt-minimal} CMake
preset flags (Ninja, \code{Release}, \code{FETCH\_ASIO=ON},
\code{MALLOC=system}), with shared linking for macOS local development.
The production \code{libraries.yaml} (Listing~2) encodes these same
flags with \code{STATIC\_LINKING=ON} for CE's sandbox.
A local \code{c++.local.properties} registered HPX with CE's frontend
(include/lib paths, \code{staticliblink} order,
\code{-DHPX\_APPLICATION\_EXPORTS}, and the macOS entry-point flag).
The production properties file (Listing~5) is derived from this PoC
configuration with Linux-specific adjustments.

% [GROUNDED: wrap/src/hpx_wrap.cpp:66-74 (macOS), 131-148 (Linux)]
\noindent\textit{Platform note:}
The macOS PoC uses \code{-Wl,-e,\_initialize\_main}; the production Linux
CE configuration (Listing~5) uses \code{-Wl,-wrap=main}.  Both invoke HPX's
entry-point wrapper (\cefile{wrap/src/hpx\_wrap.cpp}): on Linux,
\code{\_\_wrap\_main()} intercepts \code{main()} via the GNU linker's
\code{-wrap} mechanism (line~136); on macOS, \code{initialize\_main()} is
set as the binary entry point via \code{-e} (line~66).  The production CE
entry uses only the Linux flag.

% [FIX 15: Non-Obvious Issues promoted and moved before Validation Results]
\subsection*{Non-Obvious Issues Discovered During PoC}

Three non-obvious bugs surfaced that are invisible from reading
documentation alone:

\begin{enumerate}[leftmargin=2em, itemsep=2pt]
  \item \textbf{Boost headers on include path:} HPX headers transitively
        include \code{<boost/version.hpp>}.  CE must add Boost's include
        directory alongside HPX's.  The \code{\%DEP0\%} mechanism handles
        this in production; for the PoC, I added \code{/opt/homebrew/include}.

  \item \textbf{\code{hpx\_main.hpp} entry point wrapping:} HPX's
        zero-boilerplate mode requires the linker to wrap \code{main()}.
        Linux: \code{-Wl,-wrap=main}.  macOS:
        \code{-Wl,-e,\_initialize\_main}.  Without this flag:
        \code{Attempting to use hpx\_main.hpp functionality without
        linking to libhpx\_wrap}.

  % [GROUNDED: libs/CMakeLists.txt:80 (hpx_full → libhpx), :93 (hpx_full→hpx_core), :152 (hpx INTERFACE)]
  \item \textbf{Static library link order:} HPX's libraries have strict
        dependency order: \code{hpx\_wrap} $\to$ \code{hpx\_init} $\to$
        \code{hpx} $\to$ \code{hpx\_core}.  The \code{hpx} library
        (built from the \code{hpx\_full} target, \cefile{libs/CMakeLists.txt:80})
        contains the full runtime and links to \code{hpx\_core}
        (\cefile{libs/CMakeLists.txt:93}).  Reversing this order or omitting
        \code{hpx} produces undefined-reference errors.
        The \code{staticliblink} property must list them in this sequence.
\end{enumerate}

\subsection*{Validation Results}

\textbf{Step~1: Library Discovery.}
The CE library search dialog discovers HPX with its description
and project website link (Figure~\ref{fig:ce-lib}).

\begin{figure}[H]
\centering
\includegraphics[width=0.75\linewidth]{ce_library_search}
\caption{HPX appears in CE's library search, confirming the properties
  file is correctly parsed by the frontend.}
\label{fig:ce-lib}
\end{figure}

\textbf{Step~2: Compilation and Execution.}
A \code{hpx::transform\_reduce} example with \code{hpx::execution::par}
compiled with Apple Clang~17 and executed in $\sim$5.8\,s (dominated by
HPX runtime bootstrap on a single-machine dev setup).  Output:
\code{Program returned: 0} / \code{Sum of squares: 204}.

% [GROUNDED: libs/full/init_runtime/src/hpx_init.cpp:46-62,121-161 — DISTRIBUTED_RUNTIME guards]
\noindent\textit{Bootstrap time on CE:}
The 5.8\,s figure is on a macOS dev machine with shared libraries and full
distributed runtime enabled.  On CE's Linux nodes with static linking and
\code{DISTRIBUTED\_RUNTIME=OFF}, bootstrap drops significantly:
\cefile{hpx\_init.cpp} guards AGAS client initialisation, parcelport
registration (\code{init\_all\_parcelports()}, lines~121--122), and
distributed action/component tables behind
\code{\#if defined(HPX\_HAVE\_DISTRIBUTED\_RUNTIME)}
(lines~46--62, 184, 222, 237, 462).
Skipping these eliminates all network socket setup and symbol-table
population.  Expected CE bootstrap: $<$500\,ms.

\begin{figure}[H]
\centering
\includegraphics[width=\linewidth]{ce_execution}
\caption{HPX \code{transform\_reduce} compiling and executing.  Left: source
  with \code{hpx::execution::par}.  Right: output confirming the runtime
  initialises, runs the parallel algorithm, and shuts down cleanly.}
\label{fig:ce-exec}
\end{figure}

\textbf{Step~3: Assembly Inspection.}
The assembly pane shows ARM64 code with HPX runtime calls
(\code{hpx::util::hpx\_logger()}, \code{hpx\_console\_logger()}).

\begin{figure}[H]
\centering
\includegraphics[width=\linewidth]{ce_assembly}
\caption{Three-pane layout: source (left), ARM64 assembly (centre),
  execution result (right) with \code{-std=c++20 -DHPX\_APPLICATION\_EXPORTS}.}
\label{fig:ce-asm}
\end{figure}

% ── Differentiation ──────────────────────────────────────────────────────────
\subsection*{Differentiation from Prior Art}

\begin{table}[H]
\centering\small
\begin{tabularx}{\linewidth}{@{}lX@{}}
\toprule
\textbf{Prior Work} & \textbf{Comparison} \\
\midrule
\rowcolor{rowgray}
Previous HPX CE entry (c.~2023) &
  Used shared linking (breaks sandbox), hardcoded Boost~1.84 path
  (breaks on updates), no \code{skip\_compilers}.  Never produced a
  working build.  My approach fixes all six root causes with static
  linking, \code{\%DEP0\%} interpolation, and a curated compiler list. \\
Kokkos integration &
  Closest analogue: compiled parallelism library, multiple versions, static
  linking, \code{staticliblink=kokkoscore:kokkoscontainers:kokkossimd}
  (search for \code{libs.kokkos.name=kokkos} in
  \cefile{c++.amazon.properties} to locate the current entry).  My design directly
  follows Kokkos's proven patterns. \\
\rowcolor{rowgray}
CE issue tracker &
  No open issues or PRs for HPX.  The broken entry exists silently:
  users get a compilation failure with no error message suggesting the
  library is misconfigured. \\
\bottomrule
\end{tabularx}
\caption{How this proposal differs from prior integration attempts.}
\end{table}

% ── Technical Approach ────────────────────────────────────────────────────────
\subsection*{Technical Approach}

\subsubsection*{1. CE Architecture Overview}

CE has four components; this project touches all four:

\begin{enumerate}[leftmargin=2em, itemsep=2pt]
  \item \textbf{TypeScript frontend:} Monaco editor, library dropdown,
        compiler selector.  Reads metadata from \cefile{c++.*.properties}.
  \item \textbf{Node.js backend:} Express server.
        \cefile{base-compiler.ts} appends \code{-I}, \code{-L}, \code{-l}
        flags.  \cefile{options-handler.ts} resolves properties.
  \item \textbf{Execution sandbox:} \code{nsjail} containers (Section~``Sandbox
        Constraints'').
  \item \textbf{Build infrastructure:} \cefile{library\_builder.py} (1652 lines)
        reads \code{libraries.yaml}, builds per-compiler, packages with Conan,
        uploads to S3.  Hard timeout: 600\,s.
\end{enumerate}

Figure~\ref{fig:ce-arch} shows how the three PRs map onto CE's
component architecture.  The rest of this section details each layer.

\begin{figure}[H]
\centering
\includegraphics[width=\linewidth]{ce_architecture}
\caption{End-to-end integration pipeline.  Colour-coded boxes show which
  PR touches each component: \textbf{PR\,\#1} (HPX repo, blue),
  \textbf{PR\,\#2} (CE infra, orange), \textbf{PR\,\#3} (CE main, green).
  Arrows show the data flow from user code to execution output.
  \textit{(Created with Excalidraw.)}}
\label{fig:ce-arch}
\end{figure}

\subsubsection*{2. Build System: \code{libraries.yaml}}

\begin{lstlisting}[language={}, caption={Proposed \texttt{libraries.yaml} entry for HPX}]
hpx:
  build_type: cmake
  depends:
    - libraries/c++/boost_bin 1.86.0
  extra_cmake_arg:
    - -DCMAKE_BUILD_TYPE=Release
    - -DHPX_WITH_TESTS=OFF
    - -DHPX_WITH_EXAMPLES=OFF
    - -DHPX_WITH_STATIC_LINKING=ON       # Bug #2 fix
    - -DHPX_WITH_DISTRIBUTED_RUNTIME=OFF  # sandbox: no network
    - -DHPX_WITH_NETWORKING=OFF           # sandbox: no sockets
    - -DHPX_WITH_FETCH_ASIO=ON            # header-only download
    - -DHPX_WITH_MALLOC=system            # no jemalloc/tcmalloc
    - -DHPX_WITH_CXX_STANDARD=20
    - -DBOOST_ROOT=%DEP0%                 # Bug #1 fix
  lib_type: static
  make_utility: ninja                     # ~2x faster than Make
  package_install: true
  repo: STEllAR-GROUP/hpx
  skip_compilers:                         # Bug #4 fix
    - g41          # GCC 4.x--9.x (all point releases)
    - g51
    - g91
    - g101         # GCC 10.x--11.x (partial C++20, ICEs)
    - g111
    - clang400     # Clang 4.x--14.x (incomplete C++20)
    - clang1401
    # ... all point releases (g42,g52,...,g115,clang1400,...)
    # CE matches IDs exactly (not regex); full list from
    # /api/compilers, assembled in W2.
  staticliblink:
    - hpx_wrap
    - hpx_init
    - hpx                                  # libs/CMakeLists.txt:80 (hpx_full)
    - hpx_core
  target_prefix: v
  targets:
    - 2.0.0                               # Bug #3 fix
    - 1.11.0
  type: github
\end{lstlisting}

Design decisions:

\begin{itemize}[leftmargin=2em, itemsep=2pt]
  \item \textbf{\code{\%DEP0\%}:} \cefile{library\_builder.py}'s
        \code{expand\_make\_arg()} replaces this with Boost's resolved
        install path.  Automatic adaptation when CE bumps Boost.

  \item \textbf{Static linking:} Mandatory for \code{nsjail}.  Kokkos
        uses the same approach.

  \item \textbf{Ninja:} HPX has 1000+ source files.  Ninja roughly halves
        build time vs.\ Make.

  \item \textbf{\code{DISTRIBUTED\_RUNTIME=OFF}:} Disables AGAS and
        parcelports, $\sim$30\% fewer sources.  Hard requirement:
        sandbox has no network.

  \item \textbf{Link order:} \code{hpx\_wrap} $\to$ \code{hpx\_init} $\to$
        \code{hpx} $\to$ \code{hpx\_core} (validated in PoC).
        % [GROUNDED: libs/CMakeLists.txt:80 (hpx_full→libhpx), :93 (hpx_full→hpx_core)]
        \code{hpx} is the ``full'' library built from the \code{hpx\_full}
        CMake target (\cefile{libs/CMakeLists.txt:80}); it is \emph{not} an
        alias for \code{hpx\_core}.
\end{itemize}

\subsubsection*{3. Dependency Chain}

% [GROUNDED: cmake/HPX_SetupBoost.cmake, cmake/FindHwloc.cmake, cmake/FindAsio.cmake]
\begin{table}[H]
\centering\small
\begin{tabularx}{\linewidth}{@{}llX@{}}
\toprule
\textbf{Dep} & \textbf{Strategy} & \textbf{Details} \\
\midrule
\rowcolor{rowgray}
Boost $\geq$1.71 & \code{\%DEP0\%} &
  CE pre-builds Boost per compiler.  \cefile{HPX\_SetupBoost.cmake:39}
  sets min version to 1.71.  CE ships 1.86.0.  If mismatch:
  \cefile{line~73} emits \code{hpx\_error("Boost could not be found...")}. \\
Asio & FetchContent &
  \code{HPX\_WITH\_FETCH\_ASIO=ON}: header-only download via
  \cefile{HPX\_SetupAsio.cmake}.  Zero build overhead. \\
\rowcolor{rowgray}
hwloc & Fallback chain &
  \cefile{FindHwloc.cmake:12--41}: searches \code{HWLOC\_ROOT},
  \code{pkg-config}, standard paths.
  \textbf{(A)}~System hwloc.
  \textbf{(B)}~CE library entry.
  \textbf{(C)}~\code{HPX\_WITH\_HWLOC=OFF} $\to$ falls back to
  \code{std::thread::hardware\_concurrency()}.  Fine for 1--4 cores. \\
\bottomrule
\end{tabularx}
\caption{Dependency resolution with concrete fallbacks.}
\end{table}

\subsubsection*{4. Sandbox Constraints: Deep Dive}

% [GROUNDED: etc/nsjail/user-execution.cfg]
CE's execution sandbox (\cefile{etc/nsjail/user-execution.cfg}) imposes
precise limits that shaped every design decision:

\begin{table}[H]
\centering\small
\begin{tabularx}{\linewidth}{@{}llX@{}}
\toprule
\textbf{Constraint} & \textbf{nsjail Value} & \textbf{HPX Implication} \\
\midrule
\rowcolor{rowgray}
Memory & \code{cgroup\_mem\_max: 209715200} (200\,MiB)
  & HPX: $\sim$2--5\,MiB/worker (task queues, TLS, scheduler).
    4 workers: $\sim$20\,MiB.  Binary: $\sim$8--15\,MiB.
    Total: $\sim$35\,MiB of 200\,MiB. \\
Processes & \code{cgroup\_pids\_max: 14}
  & HPX spawns 1 thread/core (1--4).  With main: 5--6 PIDs of 14. \\
\rowcolor{rowgray}
CPU & \code{cgroup\_cpu\_ms\_per\_sec: 500}
  & 500\,ms CPU per wall-second.  Bootstrap: $\sim$50--200\,ms;
    1M-element algo: $<$1\,s CPU. \\
Network & No network namespace &
  \code{NETWORKING=OFF} + \code{DISTRIBUTED\_RUNTIME=OFF} are
  \textbf{hard requirements}.  AGAS naming and parcelports need
  sockets. \\
\rowcolor{rowgray}
Filesystem & \code{/tmp}: 20\,MiB tmpfs, \code{noexec} &
  HPX writes no temp files.  Static libs linked at compile time. \\
File descs & \code{rlimit\_nofile: 100} &
  HPX: $\sim$8 fds (eventfd per scheduler).  Well within 100. \\
\rowcolor{rowgray}
Time & Configurable ($\sim$10\,s default) &
  Bootstrap + 1M-element algo: $<$2\,s total. \\
\bottomrule
\end{tabularx}
\caption{nsjail constraints mapped to HPX resource usage.  Every limit
  has $\geq$3$\times$ headroom.\protect\footnotemark}
\end{table}
\footnotetext{Values from
  \cefile{compiler-explorer/compiler-explorer:etc/nsjail/user-execution.cfg}.
  Verify against current CE \code{main} before submission.}

\textbf{Why \code{DISTRIBUTED\_RUNTIME=OFF} is a hard requirement, not an
optimisation:}  HPX's distributed runtime initialises AGAS (Active Global
Address Space) and parcelport networking.  AGAS binds a TCP socket for
the naming service; parcelports use TCP or MPI for inter-node communication.
The \code{nsjail} sandbox provides no network namespace and no socket
creation.  Without \code{DISTRIBUTED\_RUNTIME=OFF}, the HPX runtime would
attempt to bind a socket, receive \code{EACCES}, and abort with
\code{hpx::detail::throw\_exception()}.

\textbf{Memory layout of a statically linked HPX binary:}
\begin{itemize}[leftmargin=2em, itemsep=1pt]
  \item \code{.text} (HPX runtime + user code): $\sim$4--8\,MiB
  \item \code{.rodata} (string tables, vtables): $\sim$1--2\,MiB
  \item \code{.data/.bss} (global state): $\sim$0.5\,MiB
  \item Per-thread stack (default 2\,MiB each, 4 threads): $\sim$8\,MiB
  \item Per-thread heap (task queues, caches): $\sim$2--5\,MiB each
  \item \textbf{Total peak:} $\sim$35\,MiB of 200\,MiB budget
\end{itemize}

\subsubsection*{5. Compiler Compatibility Matrix}

% [GROUNDED: .github/workflows/ — CI matrix]
HPX's CI tests against specific compilers.  Mapped to CE's available
compilers:

\begin{table}[H]
\centering\small
\begin{tabularx}{\linewidth}{@{}llcX@{}}
\toprule
\textbf{Compiler} & \textbf{CE Versions} & \textbf{Expected}
  & \textbf{Notes} \\
\midrule
\rowcolor{rowgray}
GCC 12   & 12.1--12.4   & \textcolor{merged}{Pass}
  & Minimum GCC with reliable C++20 \\
GCC 13   & 13.1--13.3   & \textcolor{merged}{Pass}
  & Ubuntu system compiler \\
\rowcolor{rowgray}
GCC 14   & 14.1--14.2   & \textcolor{merged}{Pass}
  & HPX CI primary GCC \\
GCC 15   & 15.1         & \textcolor{merged}{Pass}
  & Trunk builds tested \\
\rowcolor{rowgray}
Clang 15 & 15.0         & \textcolor{openpr}{Risk}
  & Borderline C++20; test in W2 \\
Clang 16 & 16.0         & \textcolor{merged}{Pass}
  & C++20 modules support begins \\
\rowcolor{rowgray}
Clang 17 & 17.0         & \textcolor{merged}{Pass}
  & PoC validated (Apple Clang 17) \\
Clang 18--20 & 18--20   & \textcolor{merged}{Pass}
  & HPX CI primary: Clang~20 \\
\rowcolor{rowgray}
GCC $\leq$11 & 4.x--11.x & \textcolor{red}{Skip}
  & No C++20; ICEs on HPX templates \\
Clang $\leq$14 & 3.x--14 & \textcolor{red}{Skip}
  & Incomplete C++20 \\
\bottomrule
\end{tabularx}
\caption{Compiler matrix derived from HPX CI.  ``Skip'' compilers go
  into \code{skip\_compilers}.  ``Risk'' compilers tested in Week~2.
  GCC~10/11 have partial C++20 support but are not tested in HPX's
  x86 CI (only GCC~11 appears in ARM cross-compile workflows);
  they are included in \code{skip\_compilers} via
  \code{gcc1[01]*}.\protect\footnotemark}
\end{table}
\footnotetext{HPX CI uses \code{stellargroup/build\_env:\{17,19\}}
  containers (Clang~20/21) for x86.  GCC~12+ is the practical minimum
  for reliable C++20.  See
  \cefile{.github/workflows/build-and-test.yml}.}

\subsubsection*{6. Execution Flow}

When a user selects HPX and clicks ``Run'':

\begin{enumerate}[leftmargin=2em, itemsep=1pt]
  \item Frontend sends code + compiler ID + library selection via POST
        to \code{/api/compiler/<id>/compile}.
  \item \cefile{options-handler.ts} resolves HPX's version and path from
        \code{c++.amazon.properties}.
  \item \cefile{base-compiler.ts} appends flags:
        \code{-I .../include -L .../lib -lhpx\_wrap -lhpx\_init
        -lhpx -lhpx\_core -pthread -DHPX\_APPLICATION\_EXPORTS -Wl,-wrap=main}
  \item Compiler produces a statically linked ELF ($\sim$8--15\,MiB).
  \item Binary executes inside \code{nsjail}; \code{hpx\_main.hpp}
        bootstraps the runtime.
  \item User code runs; output captured and displayed.
\end{enumerate}

\subsubsection*{7. Default Example}

\begin{lstlisting}[language=C++, caption={Default CE example for HPX}]
// HPX: The C++ Standard Library for Parallelism and Concurrency
// Try changing 'par' to 'seq' to compare performance!

#include <hpx/hpx_main.hpp>
#include <hpx/algorithm.hpp>
#include <hpx/numeric.hpp>
#include <hpx/experimental/sandbox.hpp>
#include <iostream>
#include <numeric>
#include <vector>

int main()
{
    hpx::experimental::sandbox::describe_environment(std::cout);

    constexpr std::size_t N = 1'000'000;
    std::vector<double> data(N);
    std::iota(data.begin(), data.end(), 1.0);

    auto report = hpx::experimental::sandbox::benchmark(
        "transform_reduce (1M doubles)",
        [&]() {
            volatile double r = hpx::transform_reduce(
                hpx::execution::seq, data.begin(), data.end(),
                0.0, std::plus<>{}, [](double x) { return x * x; });
            (void)r;
        },
        [&]() {
            volatile double r = hpx::transform_reduce(
                hpx::execution::par, data.begin(), data.end(),
                0.0, std::plus<>{}, [](double x) { return x * x; });
            (void)r;
        }
    );
    report.print(std::cout);
}
\end{lstlisting}

% [GROUNDED: libs/full/include/include/hpx/experimental/sandbox.hpp:174-178,196-218]
\textbf{API signatures used in this example:}
\begin{itemize}[leftmargin=2em, itemsep=1pt]
  \item \code{void describe\_environment(std::ostream\& os)}: detects
        topology via \code{hpx::threads::create\_topology()} and prints
        cores, PUs, NUMA domains, HPX worker count, and sandbox detection
        (\cefile{sandbox.hpp:174--178}).
  \item \code{benchmark\_report benchmark(string label, SeqFn\&\& seq,
        ParFn\&\& par, size\_t iterations=5)}: runs warmup + \code{iterations}
        passes, computes mean ms, speedup, and efficiency percentage
        (\cefile{sandbox.hpp:196--218}).
  \item \code{void benchmark\_report::print(std::ostream\& os)}:
        formatted table with verdict (``Excellent / Good / Moderate / Limited /
        No speedup'') based on efficiency thresholds
        (\cefile{sandbox.hpp:93--120}).
\end{itemize}

\textbf{Estimated output on a 2-core CE node}
(\code{cgroup\_cpu\_ms\_per\_sec: 500}, effectively 0.5 CPU):
\begin{lstlisting}[style=diff, caption={Estimated CE output from default example (2-core node)}]
=== HPX Sandbox --- Environment =========================
  Cores:           2
  PUs:             2
  NUMA domains:    1
  HPX workers:     2
  Sandbox:         yes (constrained environment detected)
========================================================

=== HPX Sandbox --- Benchmark ==========================
  Label:           "transform_reduce (1M doubles)"
  Iterations:      5
  Workers:         2
--------------------------------------------------------
  Sequential:      4.82 ms
  Parallel:        2.82 ms
--------------------------------------------------------
  Speedup:         1.71x
  Efficiency:      85.5%
--------------------------------------------------------
  Verdict:         Good scaling
========================================================
\end{lstlisting}
\noindent\textit{These values are estimates; actual numbers depend on CE's
cgroup throttling and node hardware.  They will be replaced with measured
values during W7 validation.}

\subsubsection*{8. Frontend Configuration}

\begin{lstlisting}[language={}, caption={Proposed c++.amazon.properties}]
libs.hpx.name=HPX
libs.hpx.description=The C++ Standard Library for Parallelism and Concurrency
libs.hpx.url=https://hpx.stellar-group.org
libs.hpx.versions=200:110
libs.hpx.versions.200.version=2.0.0
libs.hpx.versions.200.path=/opt/compiler-explorer/libs/hpx/v2.0.0/include
libs.hpx.versions.200.libpath=/opt/compiler-explorer/libs/hpx/v2.0.0/lib
libs.hpx.versions.200.staticliblink=hpx_wrap:hpx_init:hpx:hpx_core
libs.hpx.versions.200.liblink=pthread:dl:rt
libs.hpx.versions.200.dependencies=boost
libs.hpx.versions.200.options=-DHPX_APPLICATION_EXPORTS -Wl,-wrap=main
libs.hpx.versions.110.version=1.11.0
libs.hpx.versions.110.path=/opt/compiler-explorer/libs/hpx/v1.11.0/include
libs.hpx.versions.110.libpath=/opt/compiler-explorer/libs/hpx/v1.11.0/lib
libs.hpx.versions.110.staticliblink=hpx_wrap:hpx_init:hpx:hpx_core
libs.hpx.versions.110.liblink=pthread:dl:rt
libs.hpx.versions.110.dependencies=boost
libs.hpx.versions.110.options=-DHPX_APPLICATION_EXPORTS -Wl,-wrap=main
libs.hpx.example=hpx-sandbox
\end{lstlisting}

\noindent\textit{Note:} \code{liblink=pthread:dl:rt} adds system libraries
required by \code{libhpx\_core.a} on Linux: \code{dl} for dynamic symbol
loading (\code{dlopen}/\code{dlsym}) and \code{rt} for POSIX timers
(\code{clock\_gettime}).  These are separate from \code{staticliblink}
because they are shared system libraries, not HPX-provided archives.

\subsubsection*{9. CMake Preset}

% [GROUNDED: CMakePresets.json — godbolt-minimal preset; VERIFY line numbers before submission]
HPX already ships a \code{godbolt-minimal} preset.  Current state
(verbatim from \cefile{CMakePresets.json}, search for
\code{"name": "godbolt-minimal"}):

\begin{lstlisting}[language={}, caption={Current godbolt-minimal preset (HPX source, verbatim)}]
{
  "name": "godbolt-minimal",
  "displayName": "Godbolt / Sandboxed Environment",
  "description": "Minimal single-node static build optimized for
    browser-based compilers (Compiler Explorer / Godbolt).",
  "generator": "Ninja",
  "binaryDir": "${sourceDir}/build/godbolt-minimal",
  "cacheVariables": {
    "CMAKE_BUILD_TYPE": "Release",
    "HPX_WITH_MALLOC": "system",
    "HPX_WITH_NETWORKING": "OFF",
    "HPX_WITH_TESTS": "OFF",
    "HPX_WITH_EXAMPLES": "OFF",
    "HPX_WITH_DOCUMENTATION": "OFF",
    "HPX_WITH_STATIC_LINKING": "ON",
    "HPX_WITH_FETCH_ASIO": "ON"
  }
}
\end{lstlisting}

PR\,\#1 will add three variables.  Proposed additions:

\begin{lstlisting}[style=diff, caption={Proposed additions to godbolt-minimal (PR \#1)}]
+ "HPX_WITH_DISTRIBUTED_RUNTIME": "OFF",
+ "HPX_WITH_TOOLS": "OFF",
+ "HPX_WITH_CXX_STANDARD": "20"
\end{lstlisting}

% [GROUNDED: libs/full/init_runtime/src/hpx_init.cpp:46-62 — DISTRIBUTED_RUNTIME guards]
\textbf{\code{HPX\_WITH\_DISTRIBUTED\_RUNTIME=OFF}:}
Prevents AGAS from binding a TCP socket at init and disables parcelport
registration (\cefile{hpx\_init.cpp:121--122}, guarded by
\code{HPX\_HAVE\_DISTRIBUTED\_RUNTIME} at line~46).  Mandatory in
\code{nsjail}'s network-isolated sandbox.  Also removes all action/component
registration code, reducing both binary size and compile time.

% [GROUNDED: CMakeLists.txt:468,2714-2717 — HPX_WITH_TOOLS]
\textbf{\code{HPX\_WITH\_TOOLS=OFF}:}
Disables the \cefile{tools/} subdirectory (\cefile{CMakeLists.txt:2714--2717}),
which contains development utilities (\code{hpxdep}, \code{inspect},
\code{hpx-top.py}, performance test scripts).  These are never used on CE.

\textbf{\code{HPX\_WITH\_CXX\_STANDARD=20}:}
Required for \code{hpx::experimental::sandbox} (uses C++20 concepts and
\code{std::chrono} duration formatting) and matches the \code{-std=c++20}
flag passed to user code.

\subsubsection*{10. Risks and Mitigation}

\begin{table}[H]
\centering\small
\begin{tabularx}{\linewidth}{@{}p{2.5cm}lX@{}}
\toprule
\textbf{Risk} & \textbf{P} & \textbf{Mitigation} \\
\midrule
\rowcolor{rowgray}
Build $>$ 600\,s &
  Med &
  Ninja + \code{DISTRIBUTED\_RUNTIME=OFF} + \code{TESTS=OFF}: $\sim$30\%
  fewer sources.  Precedent: Qt/Boost have custom timeouts. \\
hwloc missing &
  Low &
  Three-level fallback: system $\to$ CE entry $\to$ \code{OFF}. \\
\rowcolor{rowgray}
CE review delay &
  Med &
  Join CE Discord during bonding.  Draft PRs in W5: 5+ weeks buffer. \\
Compiler ICEs &
  Med &
  Comprehensive \code{skip\_compilers} from W2.  Known: GCC$\leq$11 ICEs. \\
\rowcolor{rowgray}
Link order &
  Low &
  Validated in PoC: \code{wrap}$\to$\code{init}$\to$\code{hpx}$\to$\code{core}. \\
Boost mismatch &
  Low &
  \cefile{HPX\_SetupBoost.cmake} emits clear error if $<$1.71.
  CE ships 1.86.0. \\
\bottomrule
\end{tabularx}
\caption{Risk analysis. P = probability.}
\end{table}

% ── Milestones ────────────────────────────────────────────────────────────────
\subsection*{Milestones and Schedule}

\begin{longtable}{|c|p{2.2cm}|p{5.8cm}|p{3.5cm}|}
\hline
\rowcolor{rowgray}
\textbf{Wk} & \textbf{Phase} & \textbf{Task} & \textbf{Definition of Done} \\
\hline
\endfirsthead
\hline
\rowcolor{rowgray}
\textbf{Wk} & \textbf{Phase} & \textbf{Task} & \textbf{Definition of Done} \\
\hline
\endhead
\hline
\endfoot

B1--B2 & Community Bonding &
  Join CE Discord.  Study \cefile{library\_builder.py}'s
  \code{expand\_make\_arg()} and \code{build\_library()}.
  Set up local CE.  Build HPX with \code{godbolt-minimal}.
  Check hwloc on CE nodes.
  & Can trace \code{expand\_make\_arg()} flow; local CE serves requests. \\
\hline

W1 & Local Setup &
  Fork CE repos.  Add HPX \code{libraries.yaml}.
  First build with GCC~14.  Verify Conan packaging.
  & HPX builds via CE pipeline on GCC~14; Conan succeeds. \\
\hline

W2 & Multi-Compiler &
  Test GCC~\{12,13,14,15\}, Clang~\{15--20\}.
  Record times/sizes.  Build \code{skip\_compilers}.
  \newline\textit{\textbf{Highest-risk week.}}
  \textbf{Contingency:} If $>$2 fail, narrow to GCC~\{14,15\}
  + Clang~\{17,20\} and expand later.
  & Matrix complete; $\geq$6 compilers pass. \\
\hline

W3 & Dependencies &
  Resolve hwloc (A/B/C).  Validate \code{\%DEP0\%}.
  Test \code{FETCH\_ASIO} across compilers.
  & All deps resolved; no manual paths. \\
\hline

W4 & Build Pipeline &
  Finalise \code{libraries.yaml}.  Timeout mitigations.
  Validate Conan packaging.
  & Full YAML produces packages for all compilers. \\
\hline

W5 & E2E Validation &
  Full pipeline: HPX $\to$ compile $\to$ execute $\to$ output.
  \textbf{Open draft PRs.}
  & Local CE works; drafts posted on CE Discord. \\
\hline

W6 & Draft PR Feedback &
  Address draft PR feedback from CE Discord; begin properties
  file skeleton.  Refine \code{skip\_compilers} based on
  maintainer guidance.
  & At least one CE maintainer has responded to draft PR. \\
\hline

\rowcolor{rowgray}
MT & \textbf{Midterm} &
  HPX on $\geq$8 compilers; draft PRs open.
  & Midterm criteria met. \\
\hline

W7 & Frontend &
  Write \code{c++.amazon.properties}.
  Verify dropdown and flag resolution.
  & HPX in dropdown; compilation succeeds from UI. \\
\hline

W8 & Example + PR\,\#1 &
  Default example + \code{sandbox.hpp} refinements.
  Submit PR\,\#1 to HPX.
  & PR\,\#1 passes HPX CI; example runs. \\
\hline

W9 & Edge Cases &
  Test: missing \code{hpx\_main.hpp}, incompatible compilers,
  timeout, HPX+Boost conflicts.  Draft \code{CONTRIBUTING\_CE.md}.
  & All edge cases documented. \\
\hline

W10 & Upstream PRs &
  Submit PR\,\#2 (CE infra) and PR\,\#3 (CE main) with
  matrix, logs, binary sizes.
  & All 3 PRs submitted. \\
\hline

W11 & Review &
  Address feedback.  Buffer for unexpected issues.
  & Review comments addressed. \\
\hline

W12 & Final &
  Final validation.  Blog post.  Maintenance docs.
  & Blog published; docs merged. \\
\hline

\rowcolor{rowgray}
FE & \textbf{Final Eval} &
  All PRs submitted; HPX functional on CE; docs complete.
  & Final criteria met. \\
\hline
\end{longtable}

\textbf{Bonding $\to$ W1 acceleration:} The bonding period studies
\cefile{library\_builder.py}'s \code{expand\_make\_arg()} and
\code{build\_library()}.  This directly accelerates W1: understanding
these functions means I write a correct YAML entry on day one, rather
than debugging syntax misunderstandings.

\subsubsection*{Gantt Chart}

\definecolor{highrisk}{RGB}{220,80,80}

\begin{center}
\begin{ganttchart}[
  hgrid, vgrid,
  x unit=0.62cm, y unit title=0.45cm, y unit chart=0.42cm,
  title label font=\scriptsize,
  bar label font=\scriptsize,
  group label font=\scriptsize\bfseries,
  bar/.append style={fill=hpxnavy!45},
  group/.append style={fill=hpxnavy!12},
  milestone/.append style={fill=red!70, inner sep=2pt},
  title height=1,
]{1}{14}
  \gantttitle{Bond}{2}
  \gantttitlelist{W1,W2,W3,W4,W5,MT,W7,W8,W9,10,11,12}{1} \\

  \ganttgroup{Phase 1: Build Pipeline}{1}{7} \\
  \ganttbar{CE study + setup}{1}{2} \\
  \ganttbar[bar/.style={fill=highrisk!50}]{Multi-compiler (\textbf{HIGH RISK})}{3}{4} \\
  \ganttbar{Dependency resolution}{5}{5} \\
  \ganttbar{YAML + E2E pipeline}{6}{7} \\

  \ganttmilestone{Midterm Evaluation}{8} \\

  \ganttgroup{Phase 2: Frontend + PRs}{9}{14} \\
  \ganttbar{Properties + UI}{9}{9} \\
  \ganttbar{Example + HPX PR}{10}{10} \\
  \ganttbar{Edge cases + docs}{11}{11} \\
  \ganttbar{Upstream PRs}{12}{12} \\
  \ganttbar{Review iteration}{13}{13} \\
  \ganttbar{Final validation}{14}{14} \\

  \ganttmilestone{Final Evaluation}{14}
\end{ganttchart}
\end{center}

\noindent\textit{Note:} The \colorbox{highrisk!50}{red bar} marks the
multi-compiler week (W2--W3), the highest-risk task identified in the
risk analysis.  Milestones (\textbf{$\blacklozenge$}) coincide with
GSoC evaluation deadlines.

% ══════════════════════════════════════════════════════════════════════════════
\section*{Why Me}

This project requires a rare combination of skills.  I am the only applicant
who has \textbf{already built the integration} and published a working
proof of concept:

\begin{enumerate}[leftmargin=2em, itemsep=3pt]
  \item \textbf{Working proof of concept:} HPX~2.0.0 compiles and executes
        inside a local CE instance with correct assembly output
        (Figures~\ref{fig:ce-lib}--\ref{fig:ce-asm}).  Published at
        \href{https://github.com/GitMasterJatin/Integrating-HPX-with-Compiler-Explorer-}%
        {GitHub}.  The hardest unknowns are resolved.

  \item \textbf{10 merged HPX PRs} spanning heap-corruption fixes (Valgrind/
        ASan diagnostics $\to$ compiler-specific build debugging), CMake
        modernisation ($\to$ YAML and preset design), scheduler internals
        ($\to$ runtime init understanding), and API cleanup ($\to$ C++20-clean
        examples).  See Table~1 for the full skill-to-task mapping.
        Four highlights:
        \begin{itemize}[leftmargin=1.5em, itemsep=1pt, topsep=2pt]
          \item PR \#7051 (heap-corruption fix via Valgrind/ASan) $\to$
                diagnosing build failures across GCC/Clang versions in W2.
          \item PRs \#7032, \#6996, \#6977, \#6976 (CMake modernisation) $\to$
                designing \code{skip\_compilers} and per-version preset flags.
          \item PR \#7027 (NUMA topology cleanup) $\to$ hwloc dependency
                handling and \code{--hpx:bind} fallback strategy.
          \item PR \#7114 (SPSC queue refactor) $\to$ understanding
                \code{hpx\_init} bootstrap sequence and thread-pool init.
        \end{itemize}

  \item \textbf{Six root causes identified and fixed:} Not ``HPX is broken''
        but a line-by-line audit with concrete fixes validated in the PoC.

  \item \textbf{CE architecture understanding:} I have read
        \cefile{library\_builder.py} (1652 lines),
        \cefile{base-compiler.ts}, \cefile{options-handler.ts}, and
        \cefile{user-execution.cfg}.  I know how \code{\%DEP0\%} expands,
        how \code{staticliblink} becomes \code{-l} flags, and how
        \code{nsjail}'s \code{cgroup\_mem\_max} constrains execution.

  \item \textbf{Full-stack web experience:} Production Node.js/Express
        (Fiverse) and TypeScript monorepo (SyncCanvas), CE's exact stack.

  \item \textbf{Community engagement:} Active on HPX Discord and mailing list.
        Will join CE Discord during bonding.
\end{enumerate}

% ══════════════════════════════════════════════════════════════════════════════
\section*{Future Work \& Stretch Goals}

\subsubsection*{Committed post-GSoC maintenance}

\begin{description}[leftmargin=1.5em, itemsep=3pt]
  \item[Version maintenance:]
    HPX releases twice a year.  Each release: update one YAML line,
    one properties line, verify builds, submit PR.

  \item[Example gallery:]
    Futures/continuations (\code{hpx::async}, \code{.then()}), parallel
    sort, \code{hpx::dataflow}, sender/receiver (P2300).

  \item[Complex-library guide:]
    Document the integration process for other compiled libraries.
\end{description}

\subsubsection*{Stretch goals (if ahead of schedule)}

\begin{enumerate}[leftmargin=2em, itemsep=3pt]
  \item \textbf{Nightly \code{master} builds:} HEAD tracking target
        in \code{libraries.yaml} for bleeding-edge testing.
  \item \textbf{CMake project mode:} Template \code{CMakeLists.txt}
        with \code{find\_package(HPX)} for multi-file examples.
  \item \textbf{Windows / MSVC:} HPX on MSVC targets via CE's
        Wine/cross-compilation.
  \item \textbf{HPX-branded CE instance:} Dedicated instance with
        extended timeouts, curated examples, and task-graph visualisation.
\end{enumerate}

% ══════════════════════════════════════════════════════════════════════════════
\subsection*{Availability}

\begin{itemize}[leftmargin=2em, itemsep=2pt]
  \item \textbf{Time:} 15--20 hours/week (175 hours total).
  \item \textbf{Dates:} May 26 -- August 18, 2026 (12 coding weeks).
  \item \textbf{Constraints:} Semester ends early May; no exams during GSoC.
        GSoC is my primary focus.
  \item \textbf{Communication:} Daily WIP pushes + dev-log;
        weekly written update to Hartmut Kaiser and Giannis Gonidelis
        (HPX Discord or email); bi-weekly 30-min video call
        (proposed Thursdays 18:00 IST\,/\,12:30 UTC).
  \item \textbf{Community:} CE Discord \code{\#library-additions}
        (joined during bonding) + HPX Discord; milestone write-ups
        with demo recordings.
\end{itemize}

\vfill
\begin{center}
\small\itshape
PoC: \url{https://github.com/GitMasterJatin/Integrating-HPX-with-Compiler-Explorer-}\\
Source: \url{https://github.com/GitMasterJatin}
\end{center}

\end{document}
